\documentclass[a4paper,11pt]{article}
\usepackage[utf8]{inputenc}
\usepackage[T1]{fontenc}
\usepackage[french]{babel}
\usepackage[left=2cm,right=2cm,top=2cm,bottom=2cm]{geometry}
\usepackage{xcolor}
\usepackage{hyperref}
\usepackage{fontawesome5}
\usepackage{parskip}

% --- Couleurs ---
\definecolor{primary}{RGB}{35, 55, 123}

\begin{document}

% --- EN-TÊTE ---
\begin{center}
    {\Large \textbf{Loïc Aron MBASSI EWOLO}}\\[4pt]
    \small
    \faMapMarker\ Cergy, France \quad
    \faPhone\ (+33) 7 59 52 33 98 \quad
    \faEnvelope\ \href{mailto:loicaron.mbassiewolo@ensea.fr}{loicaron.mbassiewolo@ensea.fr}
\end{center}
\vspace{0.5cm}

% --- DESTINATAIRE ---
\hfill
\begin{minipage}[t]{0.45\textwidth}
    \textbf{CEA Grenoble}\\
    Service Recrutement\\
    Grenoble (38)
\end{minipage}

\vspace{1cm}

\textbf{Objet :} Candidature Stage Développement et Analyse de Performance de Micro-Programmes

\vspace{0.5cm}

Madame, Monsieur,

La programmation système et l'optimisation de code sont des domaines que je pratique régulièrement, que ce soit en C sous Linux ou en Python pour des applications de calcul. Votre offre de stage sur le développement de micro-programmes pour le calcul scientifique correspond à mon profil technique. Actuellement en double diplôme à l'\textbf{ENSEA} et à l'\textbf{École Polytechnique de Yaoundé}, je souhaite mettre mes compétences au service de vos projets.

Mon projet le plus représentatif est le \textbf{projet TAXI} : j'y ai développé en C un noyau de simulation multi-processus sous Linux, avec gestion de mémoire partagée, communication inter-processus et synchronisation par signaux. Ce travail m'a confronté aux problématiques de performance et de gestion des ressources que l'on retrouve dans le calcul scientifique.

Je participe actuellement au challenge robotique \textbf{CoHoMa}, où je travaille sur l'optimisation d'algorithmes sur Nvidia Jetson. J'y fais du profilage et de l'analyse de performance pour respecter des contraintes temps réel. J'utilise Docker pour garantir la reproductibilité des environnements, et des scripts shell pour automatiser les déploiements.

Mon passage au \textbf{Mila} comme Assistant de Recherche m'a appris à travailler en autonomie sur des problématiques techniques, avec un reporting régulier. Ma formation à l'ENSEA en systèmes embarqués me donne des bases en architecture des processeurs et en programmation bas niveau qui complètent ce profil.

Je suis disponible dès \textbf{avril 2026} pour un stage de 4 à 6 mois. Je serais ravi d'échanger sur vos problématiques de performance lors d'un entretien.

Je vous prie d'agréer, Madame, Monsieur, l'expression de mes salutations distinguées.

\vspace{1cm}

\textbf{Loïc Aron MBASSI EWOLO}\\
\textit{Élève-Ingénieur ENSEA / Polytechnique Yaoundé}

\end{document}
